\section{Riemann--Roch and Sheaf Cohomology}
\label{sec:riemann-roch}

\begin{layman}
Riemann--Roch is the classical bridge between algebra and analysis on
a compact Riemann surface. It controls the dimension of the space of
meromorphic functions with prescribed zeros and poles. Pick a
\emph{divisor} \(D\) on \(X\) — a finite signed list of points, like
``allow at most a double pole here, demand at least a simple zero
there.'' The associated \emph{Riemann--Roch space} \(L(D)\) consists
of meromorphic functions whose pole/zero behavior is bounded by
\(D\). It is a complex vector space, finite-dimensional, and its
dimension \(\ell(D)\) is what Riemann--Roch computes.

The modern statement runs:
\(\ell(D) - \ell(K - D) = \deg(D) + 1 - g(X),\)
where \(K\) is the \emph{canonical divisor} (the divisor associated
to a global meromorphic 1-form) and \(g(X)\) is the analytic genus
fixed in the previous chapter. Two pieces fall out: an exact
identity (when the right side fully determines \(\ell(D)\), e.g.\
when \(\deg(D)\) is large enough) and an inequality known as
Riemann's theorem. Together they let us produce meromorphic
functions on demand — and the existence of one specific meromorphic
function (a degree-one map to the Riemann sphere on a genus-zero
surface) drives the genus-zero classification of the next chapter.

The proof of Riemann--Roch in modern form factors through two big
ingredients. \emph{Sheaf cohomology}: to a divisor \(D\) attach an
invertible sheaf \(\mathcal{O}_X(D)\); its global sections recover
\(L(D)\) and its first cohomology \(H^1(X,\mathcal{O}_X(D))\) is
the obstruction to finding more sections. \emph{Serre duality}: a
canonical pairing reduces \(H^1(X,\mathcal{L})\) to global
sections of \(K \otimes \mathcal{L}^{-1}\), making the obstruction
explicit and finite-dimensional. The Euler characteristic of an
invertible sheaf collapses to a single closed-form expression in
the divisor degree and the genus — which is the Riemann--Roch
identity in cohomological form.

Why bother? Riemann--Roch is the engine of one of the four anti-hack
theorems (genus-zero classification: \(g(X) = 0 \iff X
\cong S^2\)). Without it, you cannot produce the degree-one
meromorphic function that classifies the sphere. The chapter packages
the entire argument as a single named ``classical input'' umbrella,
\(\mathrm{input{:}riemann\text{-}roch}\), with sheaf cohomology and
Serre duality named separately so the dependency structure is
visible to readers and to the formalisation pipeline. Mathlib does
not yet contain any of this for compact Riemann surfaces; the
project ships top-down decomposition stubs under
\code{Jacobian/HolomorphicForms/Serre/} (see the Serre-duality node
for the subtree).
\end{layman}


We package Riemann-Roch as a single classical input, with named
prerequisite sheaf/cohomology infrastructure that makes the
decomposition visible in the dep graph.

\begin{definition}[Sheaf cohomology of line bundles on \(X\) (placeholder)]
\label{def:sheaf-cohomology-rs}
For an invertible sheaf \(\mathcal{L}\) on \(X\) (equivalently a divisor
class), there are cohomology groups \(H^{0}(X,\mathcal{L})\) and
\(H^{1}(X,\mathcal{L})\), with \(H^{0}(X,\mathcal{O}_X(D))=L(D)\).
\depends{Sheaves on Riemann surfaces, sheaf cohomology, line bundles —
none packaged in Mathlib for compact RS at v4.28.0.}
\lean{JacobianChallenge.HolomorphicForms.RSSheafCohomology}
\leanok
\end{definition}

\begin{layman}
A \emph{line bundle} on the surface is a way of attaching a tiny
one-dimensional complex line to every point, varying smoothly from point
to point — like a ribbon glued onto the surface. Sometimes those ribbons
can be twisted together globally; sometimes not. The two cohomology
groups \(H^0\) and \(H^1\) are linear-algebraic gadgets that count,
respectively, how many \emph{global sections} the bundle has (smooth
choices of one point on each ribbon) and how badly the bundle is twisted
(the obstruction to having sections). For a bundle coming from a
divisor — a budget of zeros and poles — the global sections are exactly
the Riemann–Roch space we already met. The first group is harder to
picture, but the next step (Serre duality) will rewrite it as global
sections of a different bundle, making it computable. Why bother?
Riemann–Roch in its modern form is a relationship between these two
groups, and stating it cleanly requires giving them names.
\end{layman}

\begin{theorem}[Serre duality on a compact Riemann surface]
\label{thm:serre-duality-rs}
\uses{def:sheaf-cohomology-rs}
For an invertible sheaf \(\mathcal{L}\) and the canonical sheaf \(K\),
\[
  H^{1}(X,\mathcal{L})\cong H^{0}(X,K\otimes\mathcal{L}^{-1})^{\vee}.
\]
\depends{Mathlib v4.28.0 has scheme-level sheaf cohomology but not
the Riemann-surface invertible-sheaf / canonical-bundle / duality
package needed here. The project ships a top-down decomposition of
the proof tree under \code{Jacobian/HolomorphicForms/Serre/}
(${\sim}30$ files: structure / cotangent / tensor presheaves and
sheaves, Dolbeault, harmonic forms, residue map, trace map,
\(L^2\) and Hodge pairings, cup-product / Yoneda, dualizing sheaf,
non-degeneracy, line-bundle Serre, Riemann--Roch corollaries) plus
a parallel bottom-up sketch in \code{Jacobian/StageB/SerreDuality.lean}.
All sub-leaves currently sorry.}
\lean{JacobianChallenge.HolomorphicForms.serre_duality_rs}
\notready
\end{theorem}

\begin{layman}
Serre duality is a magical identity. The first cohomology group of a
line bundle — the slippery, hard-to-picture one — turns out to be the
linear dual of a much friendlier object: the global sections of the
\emph{canonical} bundle twisted by the inverse of your bundle. ``Linear
dual'' just means ``space of linear functions on,'' so the dimension
matches. In short, every difficult \(H^1\) on a compact Riemann surface
quietly is a familiar \(H^0\) of something else. The canonical bundle
itself is the bundle of holomorphic 1-forms, the same finite-dimensional
space we just learned to bound. Why bother? It converts the
two-cohomology-group statement of Riemann–Roch into a single inequality
between countable, vector-space dimensions — exactly the form most useful
for producing meromorphic functions with prescribed zeros and poles.
\end{layman}

\begin{theorem}[Euler characteristic of an invertible sheaf]
\label{thm:euler-char-line-bundle}
\uses{def:sheaf-cohomology-rs,def:divisor-degree}
For an invertible sheaf \(\mathcal{L}=\mathcal{O}_X(D)\),
\(\chi(\mathcal{L}) := h^{0}(\mathcal{L}) - h^{1}(\mathcal{L})
                     = \deg(D) + 1 - g(X).\)
\lean{JacobianChallenge.HolomorphicForms.euler_char_line_bundle}
\notready
\end{theorem}

\begin{layman}
Take the dimensions of the two cohomology groups of a line bundle and
subtract: that's the \emph{Euler characteristic} of the bundle. The
theorem says this single number is determined by an embarrassingly simple
formula — the degree of the divisor (i.e., its total weight, possibly
negative), plus 1, minus the genus of the surface. No matter how
intricately the bundle is twisted, the alternating sum of its cohomology
dimensions collapses to that line. It's the same kind of magic as
Euler's classical \(V - E + F = 2\) for polyhedra: a wild combinatorial
quantity that's secretly determined by topology. Why bother? Combined
with Serre duality, this formula is exactly Riemann–Roch — the theorem
that lets us count how many independent meromorphic functions fit any
given budget of zeros and poles.
\end{layman}

\begin{theorem}[Riemann's inequality]
\label{thm:riemann-inequality}
\uses{thm:euler-char-line-bundle}
For any line bundle \(\mathcal{L}\) on a compact Riemann surface of
genus \(g\),
\[
  h^{0}(\mathcal{L}) \ge \deg(\mathcal{L}) + 1 - g .
\]
The ``weaker half'' of Riemann-Roch (1857), preceding Roch's identification
of the \(h^{1}\) correction term: the inequality drops the non-negative
\(h^{1}(\mathcal{L})\) summand from the Euler characteristic identity.
\lean{JacobianChallenge.HolomorphicForms.riemann_inequality_h0}
\leanok
\end{theorem}
\begin{proof}\leanok\end{proof}

\begin{theorem}[Classical input: Riemann-Roch (umbrella)]
\label{input:riemann-roch}
\uses{input:divisors,thm:serre-duality-rs,thm:euler-char-line-bundle}
\lean{JacobianChallenge.Blueprint.input_riemann_roch}
\leanok
For every divisor \(D\) on \(X\),
\[
  \ell(D)-\ell(K-D)=\deg(D)+1-g(X),
\]
where \(K\) is a canonical divisor and \(\ell(D)=\dim_{\C} L(D)\).
Combines the Euler-characteristic formula
(Theorem~\ref{thm:euler-char-line-bundle}) with Serre duality
(Theorem~\ref{thm:serre-duality-rs}).
\end{theorem}

\begin{layman}
We package the Riemann–Roch theorem as a single classical input.
Combining Serre duality (which converts hard cohomology into easy
cohomology) and the Euler-characteristic formula (which equates an
alternating sum to a topological constant) yields, for every divisor,
the famous identity: the dimension of functions fitting your budget,
minus the dimension of functions fitting the complementary budget,
equals degree-of-budget plus one minus genus. In plain words, the gap
between two sibling counting questions is a single number you can read
off from degree and genus. Why bother? Riemann–Roch is the workhorse
that turns abstract divisor theory into actual constructions: produce a
meromorphic function with this pole here, that zero there, and the
formula tells you whether it exists and how many such functions are
independent. Many of the proofs further down the dependency graph are
one-line corollaries of it.
\end{layman}

\begin{proposition}[Genus zero gives a degree-one meromorphic function]
\label{prop:genus-zero-degree-one-map}
\uses{input:riemann-roch}
If \(g(X)=0\), then for every \(p\in X\) there is a meromorphic function
\(f:X\to\CPone\) with exactly one simple pole at \(p\). In particular, \(f\)
has degree one.
\lean{JacobianChallenge.HolomorphicForms.homeomorphic_sphere_of_analyticGenus_eq_zero}
\leanok
\end{proposition}

\begin{layman}
Suppose the surface has genus zero — so the dimension of holomorphic
1-forms is zero. Pick any single point on the surface. Riemann–Roch,
applied to that one-point divisor, says the budget ``functions allowed
at most a simple pole there, holomorphic everywhere else'' has dimension
exactly \(2\). Two dimensions worth of functions means at least one of
them is nonconstant, and the constraint guarantees its only pole is the
one we chose. Such a function has exactly one pole on the surface —
which, viewed as a map to the Riemann sphere, has degree one (it's
one-to-one on a generic value). Why bother? This is the forward
direction of the genus-zero classification. Any genus-zero surface
admits a degree-one map to the sphere, and degree-one maps will turn out
to be honest isomorphisms — pinning the surface down as the sphere
itself.
\end{layman}

\begin{proof}\leanok
Apply Riemann-Roch to \(D=p\). Since \(\deg(p)=1\) and \(g(X)=0\),
\[
  \ell(p)-\ell(K-p)=2.
\]
The second term is zero because \(K-p\) has negative degree in genus zero.
Thus \(L(p)\) has dimension \(2\), so it contains a nonconstant meromorphic
function with at most a simple pole at \(p\). Such a function cannot have any
other pole and therefore defines a degree-one map to \(\CPone\).
\end{proof}

